\section{Integral Selmer coordinates and marked rank-three point comparisons}
\label{sec:countcoord-v62}

Use the original frame of CVI1--4. The arguments below apply to any
frame $F$ satisfying $F.\mathrm{Ok}$ over a complete characteristic-zero
discrete valuation ring $O$ with finite residue field. The groups $G,\Gamma$
are profinite and the local index set is finite. Fix a marked $g_1$,
with $\eta(g_1)=1$ and $\Psi(g_1)^2-1\in O^\times$. Write
$c=F.c$, $g^*=cg^{-1}c$, $H=\ker\eta$, $T_b=\{v_\ell\}\cup S_b$,
and $N_v=I_v\cap\ker\eta_v\cap\ker\overline\Psi_v$ inside $D_v$.
All restrictions to $D_v$ use the specified $D_v\to\Gamma$.
The sign acts on matrices by $A\mapsto JA^{-t}J$ where $J$ reverses
the basis. Write $r_z$ for the base representation with sign $\eta$ and
matrix $B(g)=\operatorname{diag}(\Psi(g),1,\Psi(g)^{-1})$.
Index rows and columns by $0,1,2$, put $w_{pq}=q-p$ and
$\psi_q(g)=\Psi(g)^{1-q}$.

Use a universal marked representation $u$ over its complete local
Noetherian $O$-algebra $R$ and an adapted weighted presentation
$\pi:P=O[[X_i:i\in\Sigma]]\twoheadrightarrow R$ with $\Sigma$ finite.
For each variable fix the witness $g_i\in H,p_i,q_i$ with
$w_i=w_{p_iq_i}$ and
$\pi X_i=u(g_i)_{p_iq_i}-B(g_i)_{p_iq_i}$.
Adaptation also identifies the images of homogeneous series with the
intrinsic weight spaces $S_j=\mathrm{wpart}(O,u,j)$. Put
$S=\mathrm{Sring}(O,u)=\sum_j S_j$, the algebra of finite sums of
homogeneous elements. The residue field of $R$ is that of $O$.
Existence of these marked data is an input to this section.

% BEGIN VSCM1
\paragraph{VSCM1. Normalized representatives with the original local conditions.}
\label{claim:VSCM1}
Fix $0\ne a\in O$, let $M=O/(a)$ and give $M$ the discrete action
$g m=\alpha(g)m$ for $\alpha=\Psi^{-1}$.
The group $F.\mathrm{ShaO}(a)$ consists precisely of the cohomology
classes represented by continuous cocycles $z$ which vanish on $N_v$
for $v\notin T_b$ and are coboundaries on the whole of $D_v$ for
$v\in T_b$. Indeed the representatives with these properties form an
$O$-submodule, by adding and scaling their local primitives. Its image
in $H^1$ is already the span in the definition.

Put $t=\Psi(g_1)$, $D(z)=t z(g_1)-t^{-1}z(g_1^*)$,
$w=t^{-1}-t=-t^{-1}(t^2-1)$ and $\delta m(g)=(\alpha(g)-1)m$.
Since $\Psi(c)=1$, we have $\Psi(g_1^*)=t^{-1}$ and
$D(\delta m)=wm$. Consequently
\[
             n_x=z-\delta\bigl(w^{-1}D(z)\bigr)
\]
is the unique representative of the class $x$ with $D(n_x)=0$.
The construction is $O$-linear even when $a$ is a nonunit divisible
by a high power of the uniformizer. Continuity follows from the
continuity of $\Psi$ modulo powers of the maximal ideal and
$(a)=\mathfrak m_O^s$ for some $s\ge0$. On $N_v$ away from
$T_b$, the frame gives $\Psi_v=1$ integrally. On $D_v$ for
$v\in T_b$, subtracting a global coboundary subtracts its primitive
from the local primitive. Thus $n_x$ satisfies both required
conditions. Finally $n_x(c)=0$ since $c^2=1$, $\alpha(c)=1$ and
$2\in O^\times$.
% END VSCM1

% BEGIN VSCM2
\paragraph{VSCM2. A square-zero marked point with prescribed Selmer cocycle.}
\label{claim:VSCM2}
Let $C=O\oplus\epsilon M$ with $\epsilon^2=0$, using $\epsilon m$
for the element $(0,m)$. Its maximal ideal is
$\mathfrak n=\mathfrak m_O\oplus\epsilon M$, since invertibility is
detected in $O$. For $s\ge1$,
$\mathfrak n^s=\mathfrak m_O^s\oplus\epsilon\mathfrak m_O^{s-1}M$.
This gives completeness and separatedness from completeness of the
two finite $O$-modules. The algebra is Noetherian, local over $O$, and
has the same residue field, so universality applies.

Put $T=E_{10}-E_{21}$. Its entries have weight $-1$ and
$-JT^tJ=T$. Thus the adjoint action of $r_z(g)$ on $OT$ is
multiplication by $\Psi(g)^{-1}$ on either sign component. Define
$r_x$ with sign $\eta$ by
\[
       r_x(g)_{\rm mat}=(1+\epsilon n_x(g)T)B(g).
\]
The square-zero multiplication rule and the cocycle identity prove
the homomorphism law. The reduction is the prescribed one. Off $T_b$
the representation is trivial on $N_v$ because both the deformation
and the base point are trivial there. At $v\in T_b$, take
$n_x|_{D_v}=\delta m_v$. Conjugating $r_z|_{D_v}$ by
$1-\epsilon m_vT$ gives $r_x|_{D_v}$, since the associated matrix
coboundary is $-m_vT-\operatorname{Ad}(r_z)(-m_vT)$. This
conjugator reduces to the identity. It gives the required triangular
basis at $v_\ell$ and the diagonal basis at each $v\in S_b$, with
the specified inertial diagonal characters.

The matrix at $c$ is the identity. Marking imposes that entries
$(0,1)$ and $(1,0)$ of $r_x(g_1)_{\rm mat}-r_x(g_1^*)_{\rm mat}$
vanish. The first has no deformation term, and the second has
deformation term $\epsilon D(n_x)=0$. The base entries are zero.
Continuity at every $\mathfrak n$-adic level follows from local
constancy of $n_x$ and the adic continuity of $B$. This proves that
$r_x$ is a point of the same marked functor, with the explicit formula
\[
     r_x(g)_{pq}=B(g)_{pq}
              +\epsilon T_{pq}\psi_q(g)n_x(g).
\]
% END VSCM2

% BEGIN VSCM3
\paragraph{VSCM3. An integral injection annihilated by every linearized relation.}
\label{claim:VSCM3}
For every $0\ne a\in O$ the map
\[
  d_a:F.\mathrm{ShaO}(a)\longrightarrow M^\Sigma,\qquad
  d_a(x)_i=T_{p_iq_i}\psi_{q_i}(g_i)n_x(g_i)
\]
is $O$-linear, is zero outside weight $-1$, and is injective. For
every $f\in\ker\pi$ it satisfies
$\sum_i[X_i]f\,d_a(x)_i=0$ in $M$. Here $[X_i]f$ is the coefficient
of the monomial $X_i$. Indeed the universal map
$h_x:R\to C$ for VSCM2 sends $\pi X_i$ to $\epsilon d_a(x)_i$.
In finitely many variables the remainder of a series after subtracting
its constant and linear terms belongs to the algebraic ideal
$(X_iX_j:i,j\in\Sigma)$. Partition its degree-at-least-two monomials
by a chosen dividing pair and factor the corresponding product out
of each of the finitely many subseries. Hence
\[
       h_x(\pi f)=[1]f+\epsilon\sum_i[X_i]f\,d_a(x)_i.
\]
The second component gives all the kernel equations. If $d_a(x)=0$,
surjectivity of $\pi$ gives zero second component on every element
of $R$, including $u(g)_{10}$ for each $g\in\Gamma$. Formula VSCM2
at that entry gives $\Psi(g)n_x(g)=0$, so $n_x=0$ and $x=0$.
This proof also covers $M=0$ and an empty variable set. Applied to
the cotangent length test of CVI5 it supplies the same $V$ for
$\operatorname{length}_O F.\mathrm{ShaO}(a)\le V$ for every $a\ne0$,
whenever the other coordinate-count inputs produce that test.
% END VSCM3

% BEGIN VSCM4
\paragraph{VSCM4. Nonpositive weight tangent vectors at the origin.}
\label{claim:VSCM4}
Assume also $F.\mathrm{VanE}$, with the local conditions
given below. For $j\le0$, let $(d_i)\in O^\Sigma$ be supported in weight
$j$ and kill every equation $\sum_i[X_i]f\,d_i$ for $f\in\ker\pi$.
Then $d_i=0$ for every $i$.

The base point over $O$ gives an $O$-algebra homomorphism $R\to O$
sending $u$ to $r_z$ and all $\pi X_i$ to zero. The finite variable
expansion of VSCM3 shows that its composite with $\pi$ is the constant
term. Thus
\[
     f\longmapsto [1]f+\epsilon\sum_i [X_i]f\,d_i
\]
kills $\ker\pi$ and is multiplicative by the product rule. It descends
to a local map $R\to O[\epsilon]/\epsilon^2$ since its first
component is the origin, a local map (the residue maps identify both
residue fields). Pushing $u$ forward preserves the local and marked
conditions, and continuity follows by passage to all maximal-ideal
powers. Write the resulting matrix as
$(1+\epsilon\gamma(g))B(g)$. Then $\gamma$ is an adjoint cocycle
and its entries are adically continuous since the $\epsilon$ component
modulo $\mathfrak m_O^s$ is detected modulo
$(\mathfrak m_O,\epsilon)^{s+1}$. A homogeneous lift of $u(g)_{pq}$
has linear coefficients only in variables of weight $w_{pq}$.
Since $\psi_q(g)$ is a unit, $\gamma(g)$ has weight $j$.
Furthermore $\gamma(c)=0$ by normalization and $\gamma$ vanishes on
$N_v$ off $T_b$ by the unramified condition.

For each $v\in T_b$, use a triangularizing basis $h=h_0+\epsilon h_1$
congruent to $1$ modulo the maximal ideal of the dual numbers. A
diagonalizing basis at $S_b$ serves as a triangularizing basis.
At the element $x\in D_v\cap\ker\eta_v$ supplied by $F.\mathrm{Ok}$,
the diagonal entries of $B(x)$ have pairwise unit differences.
The conjugate $h_0^{-1}B(x)h_0$ is upper triangular. Its first
column equation shows that the first diagonal entry is $B(x)_{00}$
because $(h_0)_{00}$ is a unit. The two lower entries of the first
column of $h_0$ are then zero by the eigenvalue differences.
The second column gives the second diagonal entry $B(x)_{11}$ and
$(h_0)_{21}=0$ in the same way. Thus $h_0$ is upper triangular.
For $g\in\ker\eta_v$ put $Z=h_1h_0^{-1}$. The matrix
\[
   \gamma(g)B(g)+B(g)Z-ZB(g)
\]
is $h_0$ times the first-order part in the chosen triangular basis
times $h_0^{-1}$. It is upper triangular, and its diagonal on
inertia is zero because the specified inertial diagonal is fixed.
Consequently at weight zero $\gamma$ vanishes on $I_v$
($\eta_v(I_v)=1$). In negative weight at $v_\ell$, dividing each
strictly lower entry by $\psi_q(g)$ identifies $\gamma(g)$ on
$\ker\eta_v$ with the coboundary $Z_j-\operatorname{Ad}(r_z(g))Z_j$
of the weight-$j$ part of $Z$. These are exactly the local conditions
in $F.\mathrm{VanE}$.

The latter supplies $0\ne b\in O$ and $Y$ with
$b\gamma(g)=Y-\operatorname{Ad}(r_z(g))Y$ for every $g$. At $c$
this gives $Y=-JY^tJ$. The infinitesimal change of
$B(g_1)-B(g_1^*)$ under this coboundary is its commutator with $Y$,
with $Y$ on the left. Its $(0,1)$ and $(1,0)$ entries vanish by the
marking. The matrix being commuted is
$\operatorname{diag}(t-t^{-1},0,t^{-1}-t)$ with $t=\Psi(g_1)$.
Its adjacent diagonal differences are units, forcing $Y_{01}=Y_{10}=0$.
The reversal identity and invertibility of $2$ now give
$Y=\operatorname{diag}(y,0,-y)$. This is invariant under the
entire base adjoint action, so $b\gamma=0$. Torsion-freeness gives
$\gamma=0$. The original adapted entry differences therefore have zero
first-order term and $d_i=0$ as claimed. This concerns $O$-valued
solutions. VSCM3 concerns solutions over $O/(a)$, which can have
positive length.
% END VSCM4

% BEGIN VSCM5
\paragraph{VSCM5. Graded coordinates determine the torus orbit of a point.}
\label{claim:VSCM5}
Assume $S_0$ is finite over $O$. Every $S_j$ is a finite $O$-module,
and $S$ is generated as an ordinary $O$-algebra by the entries
$u(g)_{pq}$ with $g\in H$. Here is an argument using compactness
in the native coefficient category. Put $z_i=\pi X_i$. They are in
$\mathfrak m_R$ by the residual condition. In particular, truncations
of a series by total degree converge after applying $\pi$ because the tail of
degree $\ge k$ lies in $(X_i)^k$, an algebraic ideal because $\Sigma$
is finite. Let $B_0$ be the $O$-algebra in $R$ generated by weight-zero
monomials in the $z_i$. It is an $O$-submodule of $S_0$, hence is
finite as an $O$-module. It is closed in $R$ because for any finite generating
set the map $O^d\to R$ is continuous in the adic topologies, $O$ is
compact by completeness and finiteness of its residue field, and $R$
is Hausdorff. Homogeneous weight-zero series are limits of such
polynomials, so $B_0=S_0$.

In weight $j$ the set of exponent vectors in $\mathbb N^\Sigma$ of
weight $j$ has finitely many coordinatewise minimal elements by
Dickson's lemma. Each monomial exponent of this weight dominates such
a minimum, and the difference has weight zero. The corresponding
finite sum of monomial multiples of $B_0$ is an $O$-finite, closed
submodule of $R$. It contains all truncations of weight-$j$ series
and hence equals $S_j$. Thus $S=O[z_i:i\in\Sigma]$. Each $z_i$ is
a difference of an entry on $H$ and a scalar, and every matrix entry
of $u(g)$ is homogeneous of weight $w_{pq}$, by the entrywise torus
conjugation formula. This proves the assertion on generators.

A point $\xi:S\to K$ into an $O$-field produces a unique representation
$r_\xi$ with sign $\eta$ and with matrices obtained by applying $\xi$
to the entries of $u$. In fact those matrices and their inverses
already have entries in $S$. The inverse determinant has weight zero because
torus conjugation on the matrix component is ordinary conjugation
even on the negative-sign coset, since $\operatorname{diag}(s,1,s^{-1})$
is fixed by the sign involution. Then use the adjugate formula.
The group identities over $R$ hold over its subring $S$, and persist
under $\xi$. Uniqueness follows entrywise.

The homogeneous decomposition of a finite sum in $S$ is unique.
Indeed a vanishing finite sum lifts to a sum of homogeneous series
in $\ker\pi$, and the homogeneous kernel condition puts each of those
components in $\ker\pi$. For $s\in K^\times$ the grading thus defines a point
$\xi_s(x)=s^j\xi(x)$ on $x\in S_j$, by addition on the finite
homogeneous decomposition and multiplicativity of weights.
Its representation is torus conjugate to $r_\xi$. Conversely, if
$r_{\xi'}=\operatorname{Ad}(\operatorname{diag}(s,1,s^{-1}))r_\xi$,
then the two maps $\xi'$ and $\xi_s$ agree on all matrix entries and
hence on the algebra they generate. In particular distinct coordinate
orbits cannot yield representations related by this torus conjugation.
% END VSCM5

% BEGIN VSCM6
\paragraph{VSCM6. Matrix isomorphism over the sign kernel detects the same orbit.}
\label{claim:VSCM6}
Now use $K=\overline{\mathbb Q}_\ell$ with its given norm, the
faithful coefficient embedding $O\subset K$ over $\mathbb Z_\ell$
(where $O$ is module-finite over $\mathbb Z_\ell$), and two
characteristic-zero marked points $r,r'$ as defined in the source. Each,
after its own torus conjugation, comes from a marked point over a DVR
$A\subset K$ finite over $O$. Put $Z_r=r(g_1)_{\rm mat}-r(g_1^*)_{\rm mat}$
and $z=(Z_r)_{00}$, using $z'$ for $r'$. Normalization at $c$
gives $r(g^*)_{\rm mat}=Jr(g)_{\rm mat}^tJ$ for $g\in H$.
Together with the two marked zeros this yields
$Z_r=\operatorname{diag}(z,0,-z)$ and likewise for $r'$.
Indeed the anti-diagonal entries of the difference vanish and the
remaining off-diagonal entries are the negatives of their reflected
partners. These identities persist under torus conjugation.

The image of $A$ is integral over $\mathbb Z_\ell$, and its maximal
ideal elements have norm $<1$. Indeed $\ell$ has positive valuation in $A$
and is a unit times a positive power of a uniformizer, while units
and their inverses have norm at most one. Reduction to $r_z$ in the
two diagonal entries at $g_1$ and $g_1^*$ therefore gives
$|z-(t-t^{-1})|<1$ and the same bound for $z'$. All these expressions
are invariant under the correcting torus conjugations. The marked
unit has $|t-t^{-1}|=1$ and $2\in O^\times$ gives $|2|=1$.
Thus $|z|=|z'|=1$ and $|z-z'|<1$, excluding $z'=-z$.

Let $M$ intertwine the matrices on $H$ and be invertible.
Applying its equation at $g_1,g_1^*$ gives $MZ_r=Z_{r'}M$.
The nonzero eigenvalues and the excluded sign force all off-diagonal
entries of $M$ to be zero, and invertibility gives
$M=\operatorname{diag}(m_0,m_1,m_2)$ with each $m_i\ne0$.
For $g\in H$ write $A_g=r(g)_{\rm mat}$ and $A'_g=r'(g)_{\rm mat}$.
Applying the intertwining identity at both $g$ and $g^*$ gives
\[
 m_i(A_g)_{ij}=m_j(A'_g)_{ij},\qquad
 m_{2-j}(A_g)_{ij}=m_{2-i}(A'_g)_{ij}.
\]
If $m_0m_2=m_1^2$, the ratios $m_i/m_j$ agree with
$s^{j-i}$ for $s=m_0/m_1$. Otherwise subtraction after cross
multiplication forces $(A_g)_{ij}=0$ whenever exactly one of $i,j$
is the middle index. On the surviving pairs choose
$s^2=m_0/m_2$ in $K$ to get the same formula.
Hence $r'$ equals the torus conjugate of $r$ on $H$. It also agrees
at $c$ since both normalized matrix components there are $1$.
As $\eta(c)=-1$, this proves equality on $\Gamma$. In particular
the source's Type I pairwise nonisomorphism test on $H$ follows for
distinct coordinate orbits of characteristic-zero points by VSCM5.
The matrix calculation applies even in its displayed block case.
% END VSCM6

% BEGIN VSCM7
\paragraph{VSCM7. The separating trace on the same coordinates.}
\label{claim:VSCM7}
For any specified $\phi\in\Gamma$, the element
$A_\phi=\operatorname{tr}(u(\phi)_{\rm mat})-
\operatorname{tr}B(\phi)\in S_0$ has origin zero. Each diagonal
entry has weight zero, constants have weight zero, and the origin sends
$u$ to $r_z$ as in VSCM4. For every coordinate point $\xi$,
\[
    \xi(A_\phi)=\operatorname{tr}(r_\xi(\phi)_{\rm mat})
                    -\operatorname{tr}B(\phi).
\]
This follows by applying $\xi$ to the three diagonal entries,
and also holds when $\phi\notin H$. Assume the arithmetic input
of CVI5 gives the strict trace separation at all good Type I
characteristic-zero points, and that the remaining point comparisons
identify the selected two-sided coordinate points as such points.
Then $A_\phi$ is nonzero on them. In the coordinate valuation bound
with $v(x)=e(O/\mathbb Z_\ell)\,v_{\ell,1}(x)$, substituting the
displayed identity gives the right-hand side
\[
   2e(O/\mathbb Z_\ell)\sum_i v_{\ell,1}
      \big(\operatorname{tr}(r_i(\phi)_{\rm mat})-
                \operatorname{tr}B(\phi)\big).
\]
Together VSCM3, VSCM5 and VSCM6 give the native finite-coefficient
length test, the orbit and matrix-isomorphism comparison, and VSCM7
gives the trace substitution. VSCM4 gives the nonpositive integral
tangent test under $F.\mathrm{VanE}$. VSCM8--10 prove the characteristic-zero realization, the
diagonal-origin comparison and the geometric tangent consequence of
$F.\mathrm{Rigid}(r)$. Deriving Type I from the two-sided condition
and supplying the arithmetic rigidity and trace hypotheses remain
inputs to this comparison.
% END VSCM7

% BEGIN VSCM8
\paragraph{VSCM8. Integral realization after scaling the graded coordinates.}
\label{claim:VSCM8}
Fix $\xi:S\to K$ over $O$ under the hypotheses of VSCM5--6. We
construct its characteristic-zero marked point.
The map $S_0\hookrightarrow R$ reflects units. If $x\in S_0$ has
inverse in $R$, its inverse has weight zero because every pair of
test maps taking equal values on $x$ also takes equal values on
$x^{-1}$. Hence $S_0$ is local with maximal ideal
$S_0\cap\mathfrak m_R$. Every image $\xi(x)$ for $x\in S_0$
is integral over $\mathbb Z_\ell$, since $S_0$ is finite over $O$.
The preimage in $S_0$ of the maximal ideal of $\mathcal O_K$ is
maximal because its quotient is a finite-dimensional domain over
$\mathbb F_\ell$, hence a field. Therefore
$|\xi(x)|<1$ for $x\in S_0\cap\mathfrak m_R$.

Write $a_i=\xi(z_i)$. For $w_i=0$ this gives $|a_i|<1$.
For $w_i>0>w_h$, apply the same fact to
$z_i^{|w_h|}z_h^{w_i}\in S_0\cap\mathfrak m_R$.
Hence there is $s\in K^\times$ with $|b_i|<1$ for all
$b_i=s^{w_i}a_i$. For nonzero $a_i$ of positive weight one needs
$|s|<|a_i|^{-1/w_i}$, and for negative weight one needs
$|s|>|a_i|^{1/|w_i|}$. The pairwise strict inequalities leave an
open interval. Missing lower or upper restrictions give endpoints
$0$ or $+\infty$. The value group $\ell^{\mathbb Q}$ is dense
in the positive reals. Take the ring of integers $A$ of the finite
extension of $\mathbb Q_\ell$ generated by $\mathrm{Frac}(O)$, $s$,
and all the $a_i$ inside $K$. This is a complete DVR finite over
$O$, and each $b_i$ lies in $\mathfrak m_A$.

Power-series evaluation at $X_i=b_i$ defines a local map
$P\to A$. For a homogeneous series $f$ of weight $j$ its evaluation
in $K$ is $s^j\xi(\pi f)$. To justify passing from polynomials to
series, endow $S_j$ with its subspace topology from $R$. For any
finite $O$-spanning list of $S_j$ from VSCM5, the continuous
surjection $O^d\to S_j$ is a quotient map by compactness and
Hausdorffness. Thus $\xi|_{S_j}$ is continuous because on $O^d$ its
composite is a finite linear combination with coefficients in $K$.
The total-degree truncations of $f$ converge in $R$ after $\pi$,
and their evaluations at $b_i$ converge in $A$. On each monomial
the asserted equality follows from its weight. The two limits agree.

For $f\in\ker\pi$, the homogeneous kernel condition puts each of
its whole weight components in $\ker\pi$. The preceding identity
makes each component's evaluation zero. The series is the limit
of its sums over $|j|\le L$ because a monomial of total degree $d$
has $|j|\le 2d$. Thus those sums exhaust every finite-degree
truncation, and the full evaluation is zero. The map descends to a
local $R\to A$ (the maximal ideal of $R$ is generated by the images
of $\mathfrak m_O$ and the $X_i$). Push $u$ forward. Locality and
continuity preserve the congruences, marking and local bases.
The entrywise homogeneous evaluation formula identifies the image in
$K$ with torus conjugation of $r_\xi$ by $\operatorname{diag}(s,1,s^{-1})$.
This constructs the required $A$ and the integral marked point even
when $A$ has a larger residue field.
% END VSCM8

% BEGIN VSCM9
\paragraph{VSCM9. Triangular signs and the unique diagonal point.}
\label{claim:VSCM9}
The matrices of $r_\xi|_H$ are upper triangular precisely when
$\xi$ vanishes on $S_j$ for every $j<0$. They are lower triangular
precisely when it vanishes on $S_j$ for every $j>0$. One direction
in each assertion follows by evaluating each matrix entry of that
weight. For the converse, take the nonnegative weight monoid in
the upper case and the nonpositive monoid in the lower case. The
adapted variables $z_i$ outside the monoid are off-diagonal entries
on $H$ and have zero evaluation. In a series homogeneous of
weight outside the monoid every monomial contains one of these
variables. Partition the monomials by the first such variable and
divide it out in each subseries. This gives a finite sum of products
of such variables and homogeneous series, proving vanishing on the
entire weight space.

Assume also $F.\mathrm{Cl1}$, which says that each locally constant character of $H$
into an abelian group which kills the $\Gamma$-conjugates of all
inertial elements lying in $H$ is trivial. If $\xi$ vanishes on
all nonzero weights, VSCM8 gives an integral point $r$ over $A$ whose
matrices on $H$ are diagonal. For each $i$ the quotient
$\kappa_i(g)=r(g)_{ii}/\psi_i(g)$ on $H$ is a character with values
in $1+\mathfrak m_A$. For $v\in T_b$ its value is one on inertia. In fact the canonical ordered local eigenlines of the diagonal matrix at a
local marked element are the coordinate lines. A triangularizing
basis congruent to one has the same ordered flag, since the three
eigenvalues have distinct residues and their spectral projectors
split each stable flag. Thus the local inertia diagonal prescribed
in that basis is the actual diagonal. The sign is positive on
$I_v$ there. For $v\notin T_b$ and $h\in I_v\cap\ker\eta_v$, put
$N=|k^\times|$ for $k=O/\mathfrak m_O$. Then $h^N\in N_v$.
Both $r$ and the base point kill $N_v$, so $\kappa_i(h)^N=1$.
The geometric sum of the first $N$ powers of $\kappa_i(h)$ reduces
to the unit $N$, so $\kappa_i(h)=1$.

For each of these inertial elements $r$ thus agrees with the base
point. Conjugation by $H$ preserves the equality, since all matrices
on $H$ are diagonal. Conjugation by $c$ preserves it because both
representations have the same value $(1,-1)$ there. This accounts
for every conjugate appearing in $F.\mathrm{Cl1}$. Reducing
$\kappa_i$ modulo $\mathfrak m_A^n$ gives a locally constant
character by the point's adic continuity. The frame hypothesis
annihilates each such character. Separatedness of $A$ gives
$\kappa_i=1$. Thus $r=r_z$ on $H$ and at $c$, hence everywhere.
The correcting torus fixes $r_z$, so the original coordinate
representation is the base point. VSCM5 then identifies $\xi$ with
the origin on the full algebra $S$. In particular, two-sided
nonvanishing on the coordinate algebra is exactly the failure of
both triangularities on $H$. Identification of this latter condition
with Type I in the count uses further arithmetic hypotheses.
% END VSCM9

% BEGIN VSCM10
\paragraph{VSCM10. Arithmetic rigidity bounds the geometric tangent by the torus direction.}
\label{claim:VSCM10}
Under the hypotheses of VSCM8, suppose $F.\mathrm{Rigid}(r_\xi)$.
Given any $O$-linear point derivation $D:S\to K$, meaning
$D(xy)=\xi(x)D(y)+\xi(y)D(x)$, there is $b\in K$ such that
\[
                      D(x)=b j\xi(x)\qquad(x\in S_j).
\]
Indeed $\xi+\epsilon D:S\to K[\epsilon]/\epsilon^2$ is an
$O$-algebra map. Apply it to the representation over $S$ constructed
in VSCM5 and call the result $\widetilde r$. Put $r=r_\xi$ and
\[
 y(g)=D\bigl(u(g)_{\rm mat}\bigr)r(g)_{\rm mat}^{-1},
 \qquad \widetilde r(g)_{\rm mat}=(1+\epsilon y(g))r(g)_{\rm mat},
\]
applying $D$ entrywise. This is an adjoint cocycle since
$\widetilde r$ is multiplicative with the same sign. The restrictions
of $D$ and $\xi$ to each finite $S_j$ are continuous by the compact
quotient argument of VSCM8. Each coefficient of $u$ is continuous
into its weight space with the subspace topology from $R$.
Matrix inversion is continuous in $\mathrm{GL}_3(K)$, so $y$ is
continuous. On $N_v$ away from $T_b$, $u$ is trivial and $D(1)=0$.

We check that the local bases required for rigidity can also be
evaluated in dual numbers. At $v\in T_b$ take a local $x\in\ker\eta_v$
with $\Psi_v(x)^2-1$ a unit. In the complete local ring $R$ its
matrix $u(x)$ has three simple, ordered residual eigenvalues. Hensel
lifting gives pairwise separated roots, and the polynomial projectors
onto the three eigenlines give canonical columns by applying the $i$th
projector to $e_i$ and dividing by its $i$th entry, a unit. The resulting
basis $h$ congruent to one is functorial for local maps. Its first
two standard flags are the local upper triangularizing flags because the
stable flags in a triangularizing basis congruent to one are split
by the same projectors, with the ordered summands determined by
their reductions. At a diagonal place this argument applies to the
individual lines. The canonical basis therefore satisfies the actual
triangular or diagonal conditions and the prescribed inertial
diagonals. Any two test maps related by conjugation with the torus
send $h$ to matrices related by that same conjugation, by uniqueness
of the ordered columns with diagonal one. The same holds for
$h^{-1}$. Thus their $(p,q)$ entries have intrinsic weight $w_{pq}$
and lie in $S$. The local identities in this basis hold over the
subring $S$ by injection into $R$. Applying $\xi+\epsilon D$ gives
the required local bases for $\widetilde r$. This use of
$F.\mathrm{Rigid}$ retains the full inertial diagonal requirement.

Rigidity now supplies $Y\in M_3(K)$ with
$y(g)=Y-\operatorname{Ad}(r(g))Y$. The universal marking at $c$
implies $y(c)=0$ and $Y=-JY^tJ$. On $H$, the variation of a
matrix $r(g)_{\rm mat}$ is $[Y,r(g)_{\rm mat}]$. Therefore the
variation of the marked difference is $[Y,Z_r]$. The universal
marked equations give zero at its entries $(0,1),(1,0)$.
By VSCM6 and VSCM8, $Z_r=\operatorname{diag}(z,0,-z)$ with $z\ne0$.
Thus $Y_{01}=Y_{10}=0$ and the reversal identity gives
$Y=b\operatorname{diag}(1,0,-1)$. The derivative of an entry
$u(g)_{pq}$ on $H$ is now $b w_{pq}r(g)_{pq}$. On the graded
algebra $S$ the rule $x\mapsto bj\xi(x)$ on degree $j$ is itself
a point derivation by additivity of weights. It agrees with $D$ on
the generating entries of VSCM5, hence everywhere. At two-sided
points this proves the required tangent-dimension bound because any two
tangents are scalar multiples of the same Euler derivation.
% END VSCM10

% BEGIN VSCM11
\paragraph{VSCM11. The exact free graded algebra to which the length lemmas apply.}
\label{claim:VSCM11}
Write $P_j$ for all series of weight $j$, and
$A^{\mathrm{gr}}=\bigoplus_jP_j\subset P$ for finite sums of such
series. The degree-zero ring $P_0$ is a complete Noetherian local
$O$-algebra with residue field $k$. Indeed the semigroup of exponents
of weight zero is generated by its finitely many minimal nonzero
elements in the coordinatewise order. Subtraction of one such
minimal element strictly lowers total degree, giving the required
factorization. Sending finitely many formal variables to the
corresponding nonconstant monomials gives a surjection from a
power-series ring over $O$ onto $P_0$. To lift a series, choose for
each exponent a factorization, assigning its coefficient to the
resulting monomial in the new variables. As the original total
degree tends to infinity so does the number of factors, since the
chosen generators have bounded degree. Substitution is therefore
well-defined and has the specified coefficients. The target is a
quotient of a complete Noetherian local ring with its natural
residue field. If there are no nonzero weight-zero exponents,
$P_0=O$ directly.

In weight $j$ the finite set of minimal monomials, including the
constant if $j=0$, generates $P_j$ as a $P_0$-module. Partition
all monomials by a dividing minimum and factor it out of the
corresponding subseries. It follows that $A^{\mathrm{gr}}$ is
generated over $P_0$ by the original variables of weights
$\{-2,-1,1,2\}$. Its augmentation $x_0$ kills every nonzero weight.
The map to $S$ is onto in each weight. If homogeneous $f_\alpha$
generate $\ker\pi$ in $P$, they generate the kernel in
$A^{\mathrm{gr}}$ as an algebraic ideal. Indeed, in any expression
$f=\sum_\alpha h_\alpha f_\alpha$ for homogeneous $f$ of weight
$j$, replace $h_\alpha$ by its component of weight $j-w(f_\alpha)$
and then sum over the finitely many weights of the given element.
Thus quotient points, their torus actions and point derivations
annihilating the ideal descend exactly to $S$.

The smooth origin chart is the inclusion of $A^{\mathrm{gr}}$ into
$\mathrm{Frac}(O)[[X_i]]$. Each $X_i$ realizes its designated
linear coefficient. A homogeneous series with zero constant and
linear coefficients belongs to $(\ker x_0)^2$ in $A^{\mathrm{gr}}$. To see this, partition its monomials by a chosen pair of dividing variables, as
in VSCM3, leaving a homogeneous factor for each pair. Take finite
sums for general elements. Conversely elements in that square
have zero constant and linear coefficients, a property unchanged
upon saturation by nonzero $O$-scalars. Thus the origin cotangent
coordinates are exactly the linear coefficients. For each
$j=-1,-2,0$, the weight-$j$ relation vectors span their cotangent
coordinate space over $\mathrm{Frac}(O)$. Otherwise a nonzero
annihilating vector over the fraction field can be cleared of
denominators. It kills the linear term of every relation because generators
in other weights pair to zero, and each generator has zero constant
term, so multiplying by a series only multiplies its linear part by
the constant coefficient of that series. VSCM4 excludes this vector.

Now use the multiplicities of NMD7, padding by zero as in CVI5.
If the number of weight $-1$ variables is zero use VSCM3 with the
empty list and $V=0$. Otherwise the preceding construction verifies
the smoothness and cotangent conditions M1 for the free graded
algebra, and $S_0$ finite over $O$ gives M2. VSCM9 gives M3, and VSCM10
and VSCM7 give M4 and M5 at each two-sided point where the indicated
arithmetic rigidity and separating trace hypotheses hold.
The normalized valuation on $K^\times$ is $e_Ov_\ell$, since units
of $O$ have norm one and $\ell$ is a unit times the $e_O$th
power of a uniformizer. It restricts to the integer-valued
valuation of $O$. SP6 now gives the finite orbit cover and its
length bound. For SP1 map $x\in F.\mathrm{ShaO}(a)$ to
$f\mapsto\sum_i [X_i]f\,d_a(x)_i$ on $A^{\mathrm{gr}}$. This
additive map into the space of $O$-linear maps is injective by
evaluation at $X_i$ and VSCM3. Choosing integral lifts of $d_a(x)_i$
only in weight $-1$ realizes each functional as the reduction of
a point derivation at $x_0$ in that weight. VSCM3 kills the specified
equations. This gives the same finite $V$ in both estimates of CVI5.
% END VSCM11
